\section{The torus}\label{sec:torus}
All spatial norms in this section are on $\TT$, unless another domain is
shown. The spaces $\XX$, $\FF$, and $\BB_{\nu,a,T}$ are those defined in
Section~\ref{sec:preliminaries}.
\subsection{Statement of the density theorem}
\begin{theorem}[Sobolev density threshold]\label{thm:main}
Equip $\FF$ with its relative $L^1(0,\infty;H^s(\TT))$ norm topology.
\begin{enumerate}[label=(\roman*),nosep]
\item For every fixed $a\in\XX$, the set $\BB_{\nu,a,T}$ is dense if $s<1/2$.
\item For zero initial velocity,
\[
 \BB^0_{\nu,T}\text{ is dense in }\FF
 \quad\Longleftrightarrow\quad s<\tfrac12.
\]
\end{enumerate}
\end{theorem}

We first construct a blowup solution close to a given smooth solution. The
force estimates yield density below the stated threshold. A separate
small-force estimate then proves non-density at and above it.
\subsection{Localization and scaling}\label{sec:packet}
\begin{lemma}\label{lem:localization}
Let $B$ be a fixed coordinate ball whose closure lies inside a torus fundamental cube. If $z$ is smooth and supported in $B$, let $z_{\R}$ be its zero extension in that cube to $\R^3$, and let $z_{\TT}$ be its periodization. For $0<s<1$,
\begin{equation}\label{eq:localization}
 \norm{z_{\TT}}_{H^s(\TT)}
 \le C_{s,B}\bigl(\norm{z_{\R}}_2+
                   \norm{z_{\R}}_{\dot H^s(\R^3)}\bigr).
\end{equation}
The constant is uniform as the support shrinks inside $B$. At $s=0$ the $L^2$ norms are equal; at $s=1$ the corresponding equality holds for the $L^2$ gradient norms.
\end{lemma}
\begin{proof}
Sobolev norms on a compact manifold are equivalent to the norms obtained
by a fixed coordinate cover and partition of unity
\cite[Chapter~4, Sections~2--3]{taylor2011}. Choose one fixed cutoff equal
to one near $\overline B$ and supported in the coordinate chart. The resulting
local-to-global map is bounded from $H^s(\R^3)$ to $H^s(\TT)$.
Since $\|z\|_{H^s(\R^3)}\le C_s(\|z\|_2+\|z\|_{\dot H^s})$,
this gives \eqref{eq:localization}. The chart and cutoff are fixed, so the
constant is independent of a smaller support of $z$. At orders zero and
one, integration over the single copy of the support gives the stated identities.
\end{proof}

Choose a compact set $K_*$ containing $K$ and the spatial projection of
$\supp F$. Fix $x_0\in B$, where $B$ is the ball in
Lemma~\ref{lem:localization}. For sufficiently small $\eps>0$, require
\[
 2\eps^2<T,\qquad x_0+\eps K_*\subset B,
 \qquad t_\eps=T-\eps^2.
\]
We extend $F$ by zero to nonpositive source times; this extension is smooth
because its temporal support is compact in $(0,\infty)$. Using this
extension and the negative-time zero extensions of $U,P$ from
Lemma~\ref{lem:packetenergy}, define
\begin{align}
 U_\eps(x,t)&=\eps^{-1}U\left(\frac{x-x_0}{\eps},
                              \frac{t-t_\eps}{\eps^2}\right),\nonumber\\
 P_\eps(x,t)&=\eps^{-2}P\left(\frac{x-x_0}{\eps},
                              \frac{t-t_\eps}{\eps^2}\right),\label{eq:scaling}\\
 F_\eps(x,t)&=\eps^{-3}F\left(\frac{x-x_0}{\eps},
                              \frac{t-t_\eps}{\eps^2}\right).\nonumber
\end{align}
The first two fields are used for $t<T$; the force is defined for all positive times. Place the rescaled whole-space solution in $B$ and then periodize it. The torus contains a single copy of its support.

\begin{proposition}\label{prop:scaling}
The fields in \eqref{eq:scaling} solve the periodic momentum equation at viscosity $\nu$, start from zero, and have unbounded speed at $T$. A spatially constant pressure adjustment enforces the mean-zero normalization. Their norms satisfy
\begin{align}
 \norm{U_\eps}_{L^\infty(0,T;L^2)}
    &=\eps^{1/2}M,&
 \norm{\nabla U_\eps}_{L^2(0,T;L^2)}
    &=\eps^{1/2}D,\label{eq:packetEscale}\\
 \norm{F_\eps}_{L^q(0,\infty;L^p)}
    &=\eps^{\alpha(p,q)}\norm{F}_{L^q(0,\infty;L^p)},&
 \alpha(p,q)&=-3+\frac3p+\frac2q,\label{eq:packetFscale}
\end{align}
for $1\le p,q\le\infty$. For $0\le s\le1$,
\begin{equation}\label{eq:packetHs}
 \norm{F_\eps}_{L^1(0,\infty;H^s(\TT))}
 \le C_s\bigl(\eps^{1/2}+\eps^{1/2-s}\bigr).
\end{equation}
In particular the force tends to zero in $L^1_tH^s_x$ for every $s<1/2$, including negative $s$.
\end{proposition}
\begin{proof}
Writing $y=(x-x_0)/\eps$ and $\sigma=(t-t_\eps)/\eps^2$, every term of the momentum equation gains the common factor $\eps^{-3}$. Incompressibility is preserved and no viscosity rescaling occurs. The temporal extension is smooth by initial vanishing. The speed identity
$\norm{U_\eps(t)}_\infty=\eps^{-1}\norm{U(\sigma)}_\infty$
proves unboundedness as $t\uparrow T$.

At a fixed time the velocity has amplitude factor $\eps^{-1}$ and spatial volume factor $\eps^3$, giving the first identity in \eqref{eq:packetEscale}. Its gradient has amplitude factor $\eps^{-2}$. Squaring and integrating in both space and time gives
$\eps^{-4}\eps^3\eps^2=\eps$ times the original dissipation integral, proving the second identity.

For finite $p,q$, the force's spatial norm gains $\eps^{-3+3/p}$ and its time norm gains $\eps^{2/q}$. The positive-time domain includes the whole transformed support because $t_\eps>0$ and the original force vanishes at negative times. This proves \eqref{eq:packetFscale}; the essential-supremum cases follow by the same change of variables with zero reciprocal exponents.

Euclidean homogeneous Sobolev scaling gives, for $0<s<1$,
\[
 \norm{\eps^{-3}F((\,\cdot-x_0)/\eps,\sigma)}
          _{\dot H^s(\R^3)}
 =\eps^{-3/2-s}\norm{F(\cdot,\sigma)}_{\dot H^s(\R^3)}.
\]
The Fourier transform gives the squared homogeneous scaling factor
$\eps^{-6}\eps^6\eps^{-3-2s}=\eps^{-3-2s}$.
Time integration contributes $\eps^2$. The inhomogeneous $L^2$ term has order $\eps^{1/2}$. Lemma~\ref{lem:localization} therefore proves \eqref{eq:packetHs}. At $s=0$ use the mixed-norm identity, and at $s=1$ use the gradient scaling. For $s<0$, the Fourier definition gives $\norm{z}_{H^s}\le\norm{z}_2$, so the $s=0$ estimate proves convergence without any homogeneous negative-order claim.
\end{proof}

The force amplitude is $\eps^{-3}\norm{F}_\infty$, even though its integrated norm tends to zero below the critical order. The homogeneous $L^1_t\dot H^{1/2}_x$ norm is scale invariant. Section~\ref{sec:critical} supplies the regularity estimate needed at that order.

\subsection{Gluing to a given smooth solution}\label{sec:insertion}
Fix a given smooth solution $(v,\pi,g)$ smooth on $[0,T+\delta]$, where
$\delta>0$, $g\in\FF$, and $v(0)=a\in\XX$. We first make the given velocity vanish near the support of $U_\eps$. This prevents nonlinear interaction terms from introducing a singularity into the new force.

Choose $R>0$ with $K_*\subset B(0,R)$. The smooth Urysohn lemma
\cite[Corollary~1.3.4]{petersen} gives an open neighborhood $V$ of $K_*$
and smooth cutoff functions satisfying
\[
\begin{aligned}
 &\theta\in C_c^\infty(\R^3),\quad 0\le\theta\le1,\quad
   \theta|_V=1,\quad \supp\theta\subset B(0,R),\\
 &\eta\in C_c^\infty(\R),\quad 0\le\eta\le1,\quad
   \eta|_{[-1,1]}=1,\quad \supp\eta\subset(-2,2).
\end{aligned}
\]
These choices are fixed independently of $\eps$; their rescalings below
localize the modification near $x_0$ in space and near $T$ in time.

\begin{lemma}\label{lem:potential}
Let $v$ be smooth and divergence free in a spatial ball centered at $x_0$. In that ball, with $y=x-x_0$, define
\begin{equation}\label{eq:potential}
 A(x,t)=\int_0^1 r\,v(x_0+ry,t)\times y\dd r.
\end{equation}
Then $\nabla\times A=v$. With the fixed cutoffs $\theta$ and $\eta$ above, put
\begin{equation}\label{eq:cutoff}
 \theta_\eps(x)=\theta((x-x_0)/\eps),\quad
 \eta_\eps(t)=\eta((t-T)/\eps^2),\quad
 w_\eps=-\nabla\times(\eta_\eps\theta_\eps A).
\end{equation}
For sufficiently small $\eps$, $w_\eps$ is a smooth, divergence-free field supported inside the coordinate ball and in
$(T-2\eps^2,T+2\eps^2)$. During the active interval before blowup of $U_\eps$,
\begin{equation}\label{eq:bgzero}
 v+w_\eps=0
 \quad\hbox{on an open neighborhood of }\supp U_\eps(t).
\end{equation}
\end{lemma}
\begin{proof}
Differentiating the radial formula and using $\nabla\cdot v=0$ gives
$\nabla\times A=v$. Hence $w_\eps$ is divergence free, and on the common
plateau of the cutoffs it equals $-v$. For sufficiently small $\eps$ the
spatial support lies strictly inside the coordinate ball and the temporal
support lies in $(0,T+\delta)$, so extension by zero is smooth.
The rescaled building block is active only for $T-\eps^2\le t<T$,
where $\eta_\eps=1$; its support lies in the spatial plateau. This proves
\eqref{eq:bgzero} and the support assertions.
\end{proof}

\begin{lemma}\label{lem:correction}
For $w_\eps$ in Lemma~\ref{lem:potential}, define
\begin{align}
 H_\eps={}&\partial_tw_\eps-\nu\Delta w_\eps
 +(v\cdot\nabla)w_\eps+(w_\eps\cdot\nabla)v
 +(w_\eps\cdot\nabla)w_\eps.\label{eq:H}
\end{align}
The force $H_\eps$ is smooth across $T$ and extends by zero to all times outside its cutoff support. Its spatial support has volume $O(\eps^3)$, and its temporal support has length $O(\eps^2)$. For every spatial multi-index $\beta$ and integer $j\ge0$,
\begin{equation}\label{eq:derivativebounds}
 |\partial_t^j\partial_x^\beta w_\eps|
 \le C_{\beta,j}\eps^{-2j-|\beta|},\qquad
 |\partial_x^\beta H_\eps|
 \le C_\beta\eps^{-2-|\beta|}.
\end{equation}
In particular,
\begin{align}
 \norm{w_\eps}_{E_T}&\le C\eps^{3/2},\label{eq:wE}\\
 \norm{H_\eps}_{L^q(0,\infty;L^p)}
   &\le C_{p,q}\eps^{-2+3/p+2/q}
     =C_{p,q}\eps^{\alpha(p,q)+1},\label{eq:Hmixed}\\
 \norm{H_\eps}_{L^1(0,\infty;H^s)}
   &\le C_s\bigl(\eps^{3/2}+\eps^{3/2-s}\bigr)
       \quad(0\le s\le1).\label{eq:HHs}
\end{align}
All constants are independent of sufficiently small $\eps$.
\end{lemma}
\begin{proof}
Set $x=x_0+\eps z$, $t=T+\eps^2\sigma$. The radial potential has
$A(x,t)=\eps\mathcal A_\eps(z,\sigma)$, where
$\mathcal A_\eps$ and all its derivatives are uniformly bounded on the
fixed cutoff support. Thus
\[
 w_\eps(x,t)=W_\eps(z,\sigma),\qquad
 W_\eps=-\nabla_z\times(\eta(\sigma)\theta(z)\mathcal A_\eps).
\]
Substitution in \eqref{eq:H} gives $H_\eps(x,t)=\eps^{-2}Q_\eps(z,\sigma)$,
where $Q_\eps$ is uniformly smooth and supported in the same fixed cylinder:
its terms are $\partial_\sigma W_\eps$, $-\nu\Delta_zW_\eps$ and
transport terms with nonnegative powers of $\eps$ and bounded coefficients.
This proves \eqref{eq:derivativebounds}. Changing variables gives
\eqref{eq:wE}--\eqref{eq:Hmixed}; applying Sobolev scaling to $Q_\eps$
and then Lemma~\ref{lem:localization} gives \eqref{eq:HHs}.
The smooth cutoffs give the asserted zero extensions in time and space.
\end{proof}

\begin{theorem}\label{thm:insertion}
Let $(v,\pi,g)$ be a periodic solution smooth on $[0,T+\delta]$, with
$\delta>0$, $g\in\FF$ and $v(0)=a\in\XX$. Fix any nonempty coordinate ball.
For all sufficiently small $\eps>0$, there are $g_\eps\in\FF$ and a solution $u_\eps$ with
\begin{enumerate}[label=(\roman*),nosep]
\item $T_{\max}^{\nu}(a,g_\eps)=T$ and
$\limsup_{t\uparrow T}\norm{u_\eps(t)}_\infty=\infty$;
\item $u_\eps=v$ for $0\le t\le T-2\eps^2$;
\item $u_\eps-v$ divergence free and supported, for every $t<T$, in a ball of diameter $O(\eps)$ inside the chosen ball;
\item the simultaneous bounds
\end{enumerate}
\begin{align}
 \norm{u_\eps-v}_{E_T}
   &\le (M+D)\eps^{1/2}+C\eps^{3/2},\label{eq:Eclose}\\
 \norm{g_\eps-g}_{L^q_tL^p_x}
   &\le C_{p,q}\bigl(\eps^{\alpha(p,q)}
                    +\eps^{\alpha(p,q)+1}\bigr),\label{eq:Fclose}\\
 \norm{g_\eps-g}_{L^1_tH^s_x}
   &\le C_s\bigl(\eps^{1/2-s}+\eps^{3/2-s}\bigr)
                         \quad(0\le s<1/2).\label{eq:Hsclose}
\end{align}
Here force norms are taken over $(0,\infty)$, and
$\alpha(p,q)=-3+3/p+2/q$ for $1\le p,q\le\infty$.
For $s<0$, the force difference also tends to zero in $L^1_tH^s_x$.
\end{theorem}
\begin{proof}
Choose the rescaled blowup solution and the correction from Lemmas~\ref{lem:potential}--\ref{lem:correction}, with all supports inside the prescribed ball, and set
\begin{equation}\label{eq:insertion}
 u_\eps=v+w_\eps+U_\eps,\qquad
 p_\eps=\pi+P_\eps,\qquad
 g_\eps=g+H_\eps+F_\eps.
\end{equation}
Initially use the compact representative $P_\eps$; subtract its spatial mean to normalize $p_\eps$.

Let $b_\eps=v+w_\eps$. Expanding the equation for the given smooth solution and \eqref{eq:H} gives
\[
 \partial_tb_\eps+(b_\eps\cdot\nabla)b_\eps
 -\nu\Delta b_\eps+\nabla\pi=g+H_\eps.
\]
Adding the building block equation produces the equation for $u_\eps$ except for
\[
 (b_\eps\cdot\nabla)U_\eps+(U_\eps\cdot\nabla)b_\eps.
\]
Both terms are identically zero. On a neighborhood of $\supp U_\eps$ this follows from $b_\eps=0$, including its derivatives, by \eqref{eq:bgzero}. Outside that support $U_\eps$ and its derivatives vanish, including at the support boundary by smoothness. Before its starting time it vanishes everywhere. Thus \eqref{eq:NS} holds exactly for \eqref{eq:insertion} on $[0,T)$.

Each summand of the velocity is divergence free. The correction is zero for
$t\le T-2\eps^2$, and the building block is zero before $T-\eps^2$, so the initial datum and the stated earlier history are unchanged. On the active building block support, $u_\eps=U_\eps$. Hence
\[
 \norm{u_\eps(t)}_\infty\ge\norm{U_\eps(t)}_\infty
\]
and the maximum norm is unbounded as $t\uparrow T$.

The velocity is smooth on every closed interval before blowup. The force $F_\eps$ is globally smooth by its original definition, and $H_\eps$ is globally smooth by Lemma~\ref{lem:correction}. Their time supports remain compact subsets of $(0,\infty)$ when $\eps$ is sufficiently small. Therefore $g_\eps\in\FF$, independently of the absence of a classical extension of $u_\eps$ at $T$. Uniqueness in Proposition~\ref{prop:local} identifies the constructed velocity with the maximal solution for $(a,g_\eps)$ up to every $T'<T$. Its lifespan is at least $T$, and the unbounded maximum norm excludes extension beyond $T$. Thus the lifespan is exactly $T$.

Finally, the triangle inequality, Proposition~\ref{prop:scaling} and
Lemma~\ref{lem:correction} give \eqref{eq:Eclose}--\eqref{eq:Hsclose}. In
\eqref{eq:Hsclose}, the lower-order inhomogeneous terms are absorbed because
$0<\eps\le1$ and $s\ge0$. For negative $s$, use
$\norm{z}_{H^s}\le\norm{z}_2$ and the $s=0$ estimate.
\end{proof}

\subsection{Density below the critical order}\label{sec:density}
\begin{proposition}\label{prop:density}
For $a\in\XX$, $\nu>0$, $T>0$ and $s<1/2$,
\[
 \forall g\in\FF\ \forall\rho>0\ \exists f\in\FF:
 \quad \norm{f-g}_{L^1_tH^s_x}<\rho,\qquad
       T_{\max}^{\nu}(a,f)\le T.
\]
\end{proposition}
\begin{proof}
Fix $g$ and $\rho$. There are two exhaustive cases.

If $T_{\max}^{\nu}(a,g)\le T$, choose $f=g$. It already belongs to
$\BB_{\nu,a,T}$ and the norm difference is zero.

If $T_{\max}^{\nu}(a,g)>T$, choose $\delta>0$ so that
$T+\delta<T_{\max}^{\nu}(a,g)$, and use its smooth solution as the given smooth solution in Theorem~\ref{thm:insertion}. For every sufficiently small $\eps$, that theorem supplies $g_\eps\in\BB_{\nu,a,T}$ with lifespan exactly $T$.
For $0\le s<1/2$, \eqref{eq:Hsclose} tends to zero. For $s<0$, its $s=0$ case and the embedding $L^2\hookrightarrow H^s$ do so. Choose $\eps$ so that the force difference is less than $\rho$ and set $f=g_\eps$.

Both choices preserve $a$, $\nu$, and $T$, and give the required approximation.
\end{proof}

\subsection{Critical regularity and completion of the proof}\label{sec:critical}
Non-density at and above the threshold is a consequence of classical
stability of the zero solution in the critical topology. For a direct
periodic robustness theorem with $L^2_tH^{-1/2}_x$ force perturbations, see
Mar\'in-Rubio, Robinson and Sadowski~\cite[Theorem~3, author manuscript
pp.~6--7]{mrrs}. The $L^1_tH^{1/2}_x$ small-force result needed here follows
from the forced Fujita--Kato theory in Danchin~\cite[Theorem~2.3.1,
pp.~47--49]{danchin}; we record only the adaptation of its hypotheses.

\begin{proposition}\label{prop:critical}
There is $c>0$, depending only on the unit torus and the stated norm conventions, such that
\begin{equation}\label{eq:smallcritical}
 g\in\FF,\qquad
 \rho:=\norm{g}_{L^1(0,\infty;H^{1/2})}<c\nu
\end{equation}
imply $T_{\max}^{\nu}(0,g)=\infty$.
\end{proposition}
\begin{proof}
First remove the mean as in Tao~\cite[Section~3, (36)--(38)]{tao2013}.
Set $m(t)=\int_0^t\overline g(\tau)\dd\tau$ and
$b(t)=\int_0^tm(\tau)\dd\tau$. In the translated coordinates,
$v(t,x)=u(t,x+b(t))-m(t)$ solves the mean-zero periodic equation with force
$h(t,x)=g(t,x+b(t))-\overline g(t)$. Translation is an isometry and removal
of the zero mode decreases the norm, so
$\|h\|_{L^1_tH^{1/2}_x}\le\rho$ and $|m(t)|\le\rho$.

Apply the critical small-data fixed-point theorem
\cite[Theorem~2.3.1 and its proof]{danchin} with $p=r=2$.
Here $\dot B^{1/2}_{2,2}=\dot H^{1/2}$, and Minkowski's inequality gives
$\|h\|_{\widetilde L^1\dot H^{1/2}}\le\|h\|_{L^1\dot H^{1/2}}$.
The same proof applies to mean-zero periodic fields using the periodic heat
estimates in \cite[Section~2.2.4, p.~44]{danchin}; the zero mode has already
been removed. It gives a global critical solution when $\rho<c\nu$.
Smooth data retain their higher regularity by the local theory cited in
Proposition~\ref{prop:local}. Restoring the translation and bounded mean
therefore gives the required global classical solution.
\end{proof}

\begin{corollary}\label{cor:nondensity}
For every $s\ge1/2$ and $T>0$, $\BB^0_{\nu,T}$ is not dense in
$\FF$ with the relative $L^1_tH^s_x$ topology.
\end{corollary}
\begin{proof}
At $s=1/2$, the set
\[
 \{g\in\FF:\norm{g}_{L^1_tH^{1/2}_x}<c\nu\}
\]
is a nonempty relative open ball, containing zero, and is disjoint from
$\BB^0_{\nu,T}$ by Proposition~\ref{prop:critical}.
For $s\ge1/2$, the Fourier weights imply
$\norm{g(t)}_{H^{1/2}}\le\norm{g(t)}_{H^s}$.
The ball $\norm{g}_{L^1_tH^s_x}<c\nu$ is therefore also disjoint from the set of forces producing breakdown. A dense subset cannot miss a nonempty open set.
\end{proof}

\begin{proof}[Proof of Theorem~\ref{thm:main}]
Proposition~\ref{prop:density} gives density for every fixed smooth initial velocity when $s<1/2$. Corollary~\ref{cor:nondensity} gives non-density for zero initial velocity when $s\ge1/2$. These are the two assertions of the theorem.
\end{proof}
The converse has been proved for zero initial velocity. The behavior for general nonzero initial velocities at or above the critical order remains outside this classification.

\subsection{Other force norms, trajectory approximation, and data pairs}
\begin{corollary}\label{cor:mixed}
For every fixed $a\in\XX$, the set $\BB_{\nu,a,T}$ is dense in $\FF$ with its relative $L^q(0,\infty;L^p)$ topology whenever
\begin{equation}\label{eq:mixedregion}
 1\le p,q\le\infty,\qquad \frac3p+\frac2q>3.
\end{equation}
\end{corollary}
\begin{proof}
Use the same two cases as in Proposition~\ref{prop:density}. In the case of a solution smooth through the prescribed time, \eqref{eq:Fclose} tends to zero because
$\alpha(p,q)>0$ and hence $\alpha(p,q)+1>0$.
\end{proof}
The sufficient region includes $L^1_tL^2_x$ and $L^2_tL^{4/3}_x$. The other mixed Lebesgue topologies require further analysis.

\begin{corollary}\label{cor:closure}
Fix $a\in\XX$. Let $\mathcal R_{a,T}$ be the trajectories with initial velocity $a$ and forces in $\FF$ that are smooth through $T$.
Let $\mathcal S_{a,T}$ be the trajectories with the same initial velocity and forces in $\FF$ that are smooth on $[0,T)$, have finite $E_T$ norm, and have unbounded maximum velocity at $T$. Then
\begin{equation}\label{eq:closure}
 \mathcal R_{a,T}\subseteq\overline{\mathcal S_{a,T}}^{\,E_T}.
\end{equation}
For each given pair $(v,g)$ and every fixed $s<1/2$, the approximating pairs can be chosen with
\[
 (u_\eps,g_\eps)\longrightarrow(v,g)
 \quad\hbox{in }E_T\times L^1(0,\infty;H^s).
\]
\end{corollary}
\begin{proof}
A given solution smooth through $T$ is smooth through $T+\delta$ for some
$\delta>0$, by local continuation; equivalently one may take this as the meaning of smoothness through the endpoint. Apply Theorem~\ref{thm:insertion}. Its energy estimate gives convergence in $E_T$, and its force estimates give simultaneous convergence. The given smooth solution has finite $E_T$ norm and the building block and correction do also, so $u_\eps$ belongs to the specified ambient energy space.

To verify the claimed interpretation of the norm, note that
\[
 \norm{z}_{L^2(0,T;L^2)}
 \le T^{1/2}\norm{z}_{L^\infty(0,T;L^2)}.
\]
Thus $E_T$ controls both terms of $L^2_tH^1_x$, while the usual sum norm controls $E_T$ directly. This proves the equivalence used in \eqref{eq:closure}. No endpoint value or classical continuation of $u_\eps$ is implied by convergence in this norm.
\end{proof}

\begin{proposition}\label{prop:projection}
Define
\[
 \mathfrak B_{\nu,T}
 =\{(a,f)\in\XX\times\FF:T_{\max}^{\nu}(a,f)\le T\}.
\]
If $s<1/2$, this set is dense in the product of any topology on $\XX$ and the relative $L^1_tH^s_x$ topology on $\FF$. Its projection onto $\XX$ is all of $\XX$.
The family obtained while requiring $a=0$ projects instead to the singleton $\{0\}$.
\end{proposition}
\begin{proof}
For every fixed $a$, Proposition~\ref{prop:density} makes the force fiber
dense and nonempty. Every product neighborhood of $(a,g)$ therefore meets
$\mathfrak B_{\nu,T}$ with its first coordinate still equal to $a$.
This proves both density and surjectivity of the projection. Restricting
to $a=0$ gives the last assertion.
\end{proof}

The projection statement has the quantifier order
$\forall a\,\exists f$, not $\exists f\,\forall a$.
For a fixed prescribed force $f_*$, the set
\[
 \{a\in\XX:T_{\max}^{\nu}(a,f_*)\le T\}
\]
is a different object and is not classified by these force-changing constructions.
There is also a distinction between the two terminal-time assertions.
Theorem~\ref{thm:insertion} gives singularity exactly at $T$ around every given solution smooth through $T$. Proposition~\ref{prop:density} uses breakdown by $T$, because a given solution may already break down earlier and is then left unchanged. It does not prove that forces producing velocity blowup exactly at $T$
are dense near a force whose solution breaks down earlier.
\begin{remark}\label{rem:peaks}
The force convergence in Theorem~\ref{thm:insertion} is compatible with diverging pointwise amplitudes. Since the original force is nonzero,
\[
 \norm{g_\eps-g}_{L^\infty_{t,x}}
 \ge\eps^{-3}\norm{F}_\infty-C\eps^{-2}\longrightarrow\infty.
\]
Here the first term is the amplitude of $F_\eps$ and the second bounds $H_\eps$. The force $F$ is nonzero by~\eqref{eq:packetenergy}, since $U$ blows up. Hence the approximating smooth forces have unbounded amplitudes and cannot obey uniform bounds on all derivatives.
\end{remark}

\subsection{Gluing in a bounded domain}
For the following bounded-domain statement, $H^s(\Omega)$ is the space
of restrictions of $H^s(\R^3)$ distributions, with quotient norm
\begin{equation}\label{eq:restriction-norm}
 \norm{z}_{H^s(\Omega)}=\inf_{Z|_\Omega=z}\norm{Z}_{H^s(\R^3)}.
\end{equation}
This convention applies also to negative orders. At order zero it is the
usual $L^2(\Omega)$ norm.
For every fixed compact $K\Subset\Omega$ and $s\in\R$, smooth fields supported
in $K$ obey
\begin{equation}\label{eq:zero-extension}
 \norm{z}_{H^s(\Omega)}\le\norm{E_0z}_{H^s(\R^3)}
 \le C_{s,K,\Omega}\norm{z}_{H^s(\Omega)},
 \qquad \operatorname{supp}z\subset K,
\end{equation}
Here $E_0z$ denotes extension by zero. Choose $\chi\in C_c^\infty(\Omega)$ equal to one near $K$ by the
smooth Urysohn lemma~\cite{petersen}. Multiplication by a fixed smooth cutoff
is bounded on $H^s(\R^3)$ for every real $s$
\cite[Chapter~4, Section~2, (2.19)]{taylor2011}. For every extension $v$ of $z$,
$\chi v=E_0z$; taking the infimum over extensions proves the second
inequality in \eqref{eq:zero-extension}. The first follows from the quotient norm.
The constant depends on the fixed cutoff and $s$, not on a smaller support
or a shrinking scale $\eps$. Applying the pointwise bounds in time gives the
same comparison for all the mixed norms in \eqref{eq:time-norms}, provided
$K$ is fixed for all times. No bounded zero-extension operator on arbitrary
$H^s(\Omega)$ is asserted. All vector and tensor norms in these definitions
are the square root of the sum of the squared component norms.

The next corollary extends the local gluing construction to a bounded
container. Placing the modification strictly inside an interior ball leaves
the velocity unchanged near the boundary, so the no-slip condition is
preserved. The starting point is a given compatible smooth solution; the
construction changes its interior evolution while retaining its boundary values.

\begin{corollary}\label{cor:boundary}
Let $\Omega\subset\R^3$ be a bounded box or a bounded smooth domain.
Suppose that, for some $\delta>0$, $(v,\pi,g)$ is a smooth no-slip
solution on $[0,T+\delta]$, and that $g$ is smooth on
$\overline\Omega\times[0,\infty)$ with temporal support compact in
$(0,\infty)$.
Here smoothness on a closed spacetime slab means restriction of a
$C^\infty$ field from an open neighborhood of that slab; this convention
also specifies smoothness at the edges and corners of a box. We assume
$v,\pi$ have this regularity on
$\overline\Omega\times[0,T+\delta]$, and $g$ on each finite closed slab.
The equation and incompressibility hold in $\Omega$, and
$v|_{\partial\Omega}=0$. Pressure may be normalized by zero spatial mean.
Existence of this compatible smooth solution is an assumption of the corollary.
The construction in Theorem~\ref{thm:insertion} applies inside any prescribed interior ball, preserving the initial velocity and no-slip boundary values.
The energy estimates use $L^2(\Omega)$ and the force estimates use
$L^1(0,\infty;H^s(\Omega))$, with the restriction norm
\eqref{eq:restriction-norm}. For every real $s$, the domain and zero-extended
force-difference norms are comparable by \eqref{eq:zero-extension}, with a
constant independent of $\eps$; the convergence assertion remains $s<1/2$.
\end{corollary}
\begin{proof}
The vector potential and all cutoff estimates are local. Choose the correction and building block supports in a smaller closed ball strictly inside $\Omega$.
The computation of the momentum equation in Theorem~\ref{thm:insertion} is unchanged. The new velocity equals the given smooth solution in a fixed boundary collar for all times before blowup, so the no-slip boundary values are preserved. The force correction is supported in the same interior region and has the same globally smooth time extension.

For $0\le s<1/2$, the Euclidean scaling proof directly bounds the zero
extensions of the force differences in $L^1(0,\infty;H^s(\R^3))$.
For $s<0$, use $\norm{E_0(g_\eps-g)}_{H^s(\R^3)}
\le\norm{E_0(g_\eps-g)}_{L^2(\R^3)}$ at each time.
All these force-difference supports, including those after $T$, lie in one
fixed compact interior ball $K$ for every sufficiently small $\eps$.
Consequently \eqref{eq:zero-extension} gives the asserted domain estimates
and the scale-independent comparison after time integration. The energy estimates follow by integrating over that ball. Finally, the classical difference-energy estimate and Gr\"onwall give uniqueness on the no-slip domain; the boundary terms vanish. Thus the constructed solution is singular exactly at $T$.
\end{proof}

\subsection{Further constructions}\label{sec:variants}
\begin{proposition}\label{prop:affine}
Fix the building block $(U,P,F)$ of Theorem~\ref{thm:packet}, and fix a spacetime cylinder $Q=B_0\times(\tau_0,\tau_1)$ with
$0<\tau_0<\tau_1<1$ and $B_0\Subset\R^3$ an open ball.
For any $b\in C_c^\infty(Q;\R^3)$ with $\nabla\cdot b=0$, define
\begin{align}
 \widetilde U&=U+b,\qquad \widetilde P=P,\nonumber\\
 \widetilde F&=F+\partial_tb-\nu\Delta b
       +(U\cdot\nabla)b+(b\cdot\nabla)U+(b\cdot\nabla)b.\label{eq:affine}
\end{align}
These fields give an infinite-dimensional affine family of distinct velocities with zero initial data, finite energy and dissipation, and the same late singular behavior as $U$. For each fixed nonzero $b$, the family obtained by replacing $b$ with $\lambda b$ is non-isolated as $\lambda\to0$ in every fixed smooth seminorm of the velocity and force differences on their compact support.
\end{proposition}
\begin{proof}
Expanding the momentum equation for $U+b$ gives exactly the force in
\eqref{eq:affine}; its divergence remains zero. Every correction term is supported in $\supp b$, a compact subset of $Q$. The original blowup solution is smooth with bounded derivatives of every fixed order on a neighborhood of that set. Thus the force correction extends smoothly by zero to the whole positive time axis, including time one, despite the singularity of $U$ there.

Since $b=0$ near time zero and near time one,
$\widetilde U(0)=0$ and $\widetilde U=U$ on a fixed late interval.
The same unbounded-speed limit holds. The energy and dissipation norms of $b$ are finite by compact smooth support, so the triangle inequality gives the corresponding bounds for $\widetilde U$.

To verify infinite dimensionality explicitly, choose countably many disjoint small balls inside $B_0$. In each ball choose a smooth compactly supported vector potential with nonzero curl, and multiply that curl by a fixed nonzero time bump in $(\tau_0,\tau_1)$. These divergence-free fields are linearly independent because their spatial supports are disjoint. Their finite linear combinations belong to the allowed class. The map $b\mapsto U+b$ is injective and affine, so its image is infinite dimensional.

For the non-isolation statement, the velocity difference is $\lambda b$, while the force difference has the form
\[
 \lambda L_U b+\lambda^2(b\cdot\nabla)b,\qquad
 L_Ub=\partial_tb-\nu\Delta b+(U\cdot\nabla)b+(b\cdot\nabla)U.
\]
On the fixed compact support, all coefficients and their derivatives are bounded. Hence for each integer $m\ge0$ the $C^m$ norm of the force difference is at most $C_m|\lambda|+C_m'\lambda^2$, and the velocity difference is at most $C_m''|\lambda|$. Both tend to zero. The forces vary in this family; no open singularity basin for one fixed force follows.
\end{proof}

\begin{proposition}\label{prop:multiple}
Fix finitely many disjoint interior balls $B_1,\ldots,B_N$ in the torus or in a bounded three-dimensional domain. At fixed $\nu>0$ and $T>0$, there is a smooth forced solution with zero initial velocity, smooth on $[0,T)$, with finite energy and dissipation, such that
\[
 \limsup_{t\uparrow T}\norm{u(t)}_{L^\infty(B_j)}=\infty
 \qquad(1\le j\le N).
\]
In a bounded domain its boundary condition is homogeneous no-slip.
\end{proposition}
\begin{proof}
Choose $x_j\in B_j$ and scales $\eps_j>0$ sufficiently small that the complete scaled compact set $x_j+\eps_jK_*$ lies strictly inside $B_j$ and $\eps_j^2<T$. Apply \eqref{eq:scaling} to each building block with start time $T-\eps_j^2$, and denote the resulting fields by $U_j,P_j,F_j$.
Set
\[
 u=\sum_{j=1}^N U_j,\qquad p=\sum_{j=1}^N P_j,\qquad f=\sum_{j=1}^N F_j.
\]
For $i\ne j$, the supports of $U_i$ and $U_j$ are separated, so
$(U_i\cdot\nabla)U_j=0$. The sum therefore satisfies the momentum equation exactly. All force supports are compact in positive time and all forces are smooth. Initial vanishing gives $u(0)=0$.

On $B_j$ the velocity equals $U_j$, proving the stated separate limsup for each ball. Disjoint supports also give
\[
 \sup_{t<T}\norm{u(t)}_2^2\le M^2\sum_{j=1}^N\eps_j,\qquad
 \int_0^T\norm{\nabla u(t)}_2^2\dd t=D^2\sum_{j=1}^N\eps_j.
\]
These sums are finite. In a bounded domain, all fields vanish in a boundary collar; hence $u=0$ on the boundary. On the torus a spatially constant pressure adjustment gives the chosen gauge.
\end{proof}
Only finitely many solutions are superposed, all with terminal time $T$. The sequences along which their maxima diverge may differ; the conclusion is a separate limsup statement in each ball.

\begin{proposition}\label{prop:conservative}
On the torus, let $f=-\nabla\phi$ for a globally defined periodic scalar
$\phi$. On a bounded domain, let $f=-\nabla\phi$ and impose homogeneous no-slip velocity. In either case, any smooth solution with zero initial velocity is identically zero on its classical lifespan.
\end{proposition}
\begin{proof}
Incompressibility and periodicity or the zero normal boundary trace give
\[
 \int f\cdot u
 =-\int\nabla\phi\cdot u=0.
\]
The classical energy identity, integrated from zero, becomes
\[
 \frac12\norm{u(t)}_2^2
 +\nu\int_0^t\norm{\nabla u(s)}_2^2\dd s=0.
\]
Both terms are nonnegative, so $u(t)=0$ at every time before blowup.
The periodic-potential condition excludes affine potentials, whose gradients need not be removable by the allowed periodic pressure gauge.
\end{proof}

Thus a nontrivial solution with zero initial velocity requires a nonconservative force.
